\section{高维数据的低维结构检验}
\label{sec:14-Deep-Learning/10-Interpretability-Mathematics/03}

你在报告里写下``图像数据位于一个低维流形附近，所以网络学得动''，配了一张用降维方法画出的二维图，上面有几团分得很开的簇。审稿人回了一句：把同样数量的独立高斯噪声点喂进同一个可视化流程，也会得到分得很开的簇，请给出别的证据。

这句反驳很难反驳。近邻类可视化方法的目标函数以保持局部邻接关系为主，全局距离和簇的分离程度会随困惑度、学习率、初始化而变，纯噪声输入照样能产生结构分明的图。于是``数据是低维的''这个在深度学习里出现频率最高的假设之一，长期停在事后解释的位置上——学得好是因为低维，低维的证据是学得好。

要脱离这个循环，得先把这句话变成可以被数据否定的陈述：在什么尺度上、用什么度量、有效维数是多少，以及这个估计对哪些处理不敏感。本节讲怎么做这件事，也讲做不到的时候该怎么讲话。

\subsection{随机投影能压缩，但压缩不是低维的证据}

先扫掉一个常见的混淆。

\begin{theorem}[Johnson--Lindenstrauss 引理的推论]
对欧氏空间中任意 \(n\) 个点，存在到 \(k=O(\varepsilon^{-2}\log n)\) 维的线性映射，使全部点对距离都落在原距离的 \(1\pm\varepsilon\) 倍之内；随机 Gaussian 投影或适当的稀疏投影以高概率给出这样的映射。
\end{theorem}

留意目标维数依赖什么：点数 \(n\) 和允许的失真 \(\varepsilon\)，与原始维数无关，也与这些点是否来自低维流形无关。均匀分布在 784 维立方体上的一万个点同样适用。它断言的是任意有限距离表都能塞进对数级的维度里，而不是数据本身有几个自由度。

因此，把数据随机投影到 \(k\) 维之后模型表现没掉，只说明当前任务对这一次压缩是稳健的。它既不识别生成数据的潜在因素，也不保证新样本、概率密度或拓扑结构被保留。用它论证内在维数，等于用一把只测长度的尺子去回答有几个自由度。

\subsection{内在维数是尺度的函数，不是数据的属性}

假设数据确实落在一个 \(d\) 维光滑流形附近。局部看，足够小的邻域内变化应当主要集中在约 \(d\) 个切方向上，邻域协方差矩阵的谱会在第 \(d\) 个位置之后陡降。这条推理成立，麻烦在于``足够小''没有绝对标准。

\begin{center}
\begin{tikzpicture}[x=1.15cm,y=0.34cm,>=stealth]
\fill[gray!15] (-1.0,0) rectangle (0.5,10.6);
\draw[->] (-3.3,0) -- (1.9,0) node[right,font=\small] {\(\log r\)};
\draw[->] (-3.3,0) -- (-3.3,11.6) node[above,font=\small] {局部有效维数};
\foreach \y in {0,2,4,6,8,10} {
  \draw (-3.42,\y) -- (-3.3,\y);
  \node[left,font=\scriptsize] at (-3.46,\y) {\y};
}
\draw[very thick] plot[smooth] coordinates
  {(-3.0,9.6) (-2.5,9.1) (-2.0,7.4) (-1.5,4.2) (-1.0,2.35)
   (-0.5,2.05) (0.0,2.0) (0.5,2.15) (1.0,3.4) (1.5,5.3)};
\draw[dashed] (-3.3,2) -- (1.6,2);
\node[right,font=\scriptsize] at (1.0,1.2) {\(d=2\)};
\node[above,font=\scriptsize,align=center] at (-0.25,10.8) {平台：可报告的估计};
\node[below,font=\scriptsize,align=center] at (-2.5,-0.5) {噪声主导\\ 看起来满维};
\node[below,font=\scriptsize,align=center] at (1.2,-0.5) {曲率与分支\\ 被碾平};
\end{tikzpicture}
\end{center}

\emph{维数估计随邻域半径变化，中间那段平台的存在与宽度本身就是证据的主要部分。}

图里的曲线来自一个可以手工构造的情形：信号躺在二维平面上，每个坐标另加一个小幅独立噪声。邻域半径远小于噪声尺度时，最近邻之间的差异几乎全是噪声，各个方向的方差相当，局部 PCA 报出接近环境维数的结果；半径远大于流形的曲率半径时，弯折、扭结与不同分支被一起碾进同一个邻域，估计又抬起来。只有中间那一段两个信号方向占主导，估计稳定在 2 附近。

这张图的读法是关键。问``数据到底几维''没有脱离尺度的唯一答案，能报告的只有``在半径 \(r\in[r_1,r_2]\) 这段区间上，估计稳定在 \(d\) 附近''。平台不出现，就说明在当前采样密度与噪声水平下，数据不支持一个尺度无关的低维描述——这本身是一个结论，不是实验失败。流形作为局部欧氏对象的定义见\hyperref[sec:12-Real-Analysis-Support/08-Manifold-Basics/01]{``流形是局部像欧氏空间的集合''}，把这个假设用于降维的具体方法见\hyperref[sec:13-Machine-Learning/04-Dimensionality-Reduction/04]{``流形假设与非线性降维''}。

\subsection{每种估计量都带着自己的一套假设}

局部 PCA 是最直观的一种：取每个点的 \(k\) 近邻，做协方差的谱分解，看特征值在哪里陡降。它的隐患在于邻域内样本数与维数往往可比——\(k=50\) 个点估计 784 维的协方差，样本协方差的谱会被系统性地拉宽，小特征值被压低、大特征值被抬高，陡降的位置随之偏移。这种失真有精确的刻画，见\hyperref[ch:094]{``随机矩阵理论：高维协方差的谱失真''}。忽略它会把维数估低或估高，方向取决于阈值怎么定。

最近邻尺度律走另一条路：若数据近似 \(d\) 维，落在半径 \(r\) 球内的样本数应当按 \(r^d\) 增长，于是把计数对半径在双对数坐标下作图，斜率就是维数。基于两个最近邻距离之比的估计量属于同一族，好处是不需要选阈值，代价是对噪声和非均匀采样敏感。相关维数与参与率则把整张距离表或整个谱压成一个标量摘要，读起来方便，但压缩过程丢掉了尺度依赖这一最重要的信息。

无论用哪一种，稳健性检查的分量不比点估计轻：换邻域大小、换距离度量、换预处理（是否标准化、是否去均值）、对样本做自助重采样，看估计漂多少。全部结论都应当以``在这些设定下''开头。

\subsection{低维不蕴含好学，好学也不蕴含低维}

一条在高维空间里绕得极其复杂的一维曲线，内在维数是 1，覆盖它的几何却可能需要海量样本，沿曲线定义的分类边界也可能高频振荡。反过来，一个支撑并不低维的问题，可能因为对称性、局部性或组合结构而完全可学——图像上的平移不变性没有降低数据的内在维数，却大幅降低了需要搜索的函数空间。

所以低维流形假设既不是深度学习成功的充分条件，也不是必要条件。把它当作唯一解释，会掩盖实际起作用的那些结构先验。

\subsection{任务需要的维数通常远小于数据的维数}

还有一层区分容易被跳过。预测目标可能只依赖数据变化中的少数方向，而整个数据分布的支撑并不低维。这两个维数没有必然联系，检验方式也不同：数据维数用前面那套几何估计，任务维数要靠干预——受控地删除或置换某组特征、做充分降维、在多个环境下比较预测的稳定性、或者直接在表示上施加扰动看输出变化。只看无监督重构保留了多少方差，可能捡回一堆方差大而与任务无关的因素，比如光照和背景。

网络也不会因为深就自动找到``真实''的生成因素。监督目标会压掉与标签无关但现实中重要的信息，也可能保留一条只在训练环境里有效的捷径。表示低维、可迁移、且对应真实机制，是三个独立的断言，需要分别给证据。

\subsection{证据不足时怎样把话说准}

多种估计量在一段合理的尺度区间内给出一致的结果，低维近似在留出数据上保住了距离、重构或预测质量，并且结论对度量选择与噪声处理不敏感——三条都满足时，你有理由说数据支持一个低维近似，并且要把尺度区间和度量一并写出来。

三条不全满足时，可以说的仍然不少：可以报告某个具体的压缩在当前任务上无损，可以报告估计随尺度如何变化而没有平台，可以报告不同估计量给出的区间彼此冲突。不能说的是``数据位于一个 \(d\) 维流形上''这种脱离尺度、脱离度量的断言。把``我们观察到的''和``我们据此相信的''分开写，是这类结论唯一稳妥的表达方式。

低维结构的检验和下一节的模型归因遇到的是同一种困难：先精确定义要解释的对象，再说清估计对哪些扰动稳定。可视化和单一分数只能提供线索，把线索当成检验，是这个领域里最容易犯也最难被发现的错误。
